% !TeX root = ../../arithmetic-and-computation.tex
\section{Detailed arithmetic underlying the $k=11$ certificate}\label{app:arithdetail}
This supplement contains the derivational arithmetic suppressed from the main proof path.  It also records the broader $42$-trace computation from which the $k=11$ point was isolated.

\section{The prime $43$: an explicit totally positive generator}\label{app:prime43}
We now specialize to the first completely split rational prime, $43$. Consider
\[
 \mathfrak l=(43,t-5)\subset\OL.
\]
Define
\begin{equation}\label{eq:alpha0}
 \alpha_0=1+\zeta_{42}+\zeta_{42}^4\in\OM.
\end{equation}
A direct resultant computation gives
\begin{equation}\label{eq:absnormalpha}
 N_{M/\Ql}(\alpha_0)=43.
\end{equation}
Its relative norm is
\begin{align}
 \pi&=\alpha_0\bar\alpha_0\\
 &=3+(\zeta+\zeta^{-1})+(\zeta^3+\zeta^{-3})+(\zeta^4+\zeta^{-4})\\
 &=t^4+t^3-4t^2-2t+5.\label{eq:pi}
\end{align}
The resultant of $f(T)$ in (\externalref{packet}{eq:minpoly}) with the polynomial in \eqref{eq:pi} is $43$, hence
\[
 N_{L/\Ql}(\pi)=43.
\]
Moreover $\pi$ is totally positive because for every real embedding $\sigma$ and either extension $\widetilde\sigma:M\hookrightarrow\C$,
\[
 \sigma(\pi)=|\widetilde\sigma(\alpha_0)|^2>0.
\]
Finally $\pi(5)\equiv0\pmod{43}$, so $(\pi)=\mathfrak l$.

\begin{proposition}\label{prop:pi1}
The element
\[
 \pi_1=t^4-t^2+t-3
\]
listed in \cite[Sec.~5.3]{Mostaed2026} has norm $43$ but is not totally positive.
\end{proposition}
\begin{proof}
At the embeddings indexed by $k=11$ and $13$ one obtains, respectively,
\[
 \sigma_{11}(\pi_1)\approx-3.1712995,\qquad
 \sigma_{13}(\pi_1)\approx-3.9795331.
\]
The signs can be certified with rational isolating intervals for the roots of (\externalref{packet}{eq:minpoly}).
\end{proof}

\begin{remark}
If $M_2^+(\OL)$ is intended in its standard Hilbert-modular sense of matrices with totally positive determinant, then the representatives $\operatorname{diag}(\pi_1,1)$ cannot lie in $M_2^+(\OL)$ as written. The totally positive generator \eqref{eq:pi} repairs this issue for the chosen prime.
\end{remark}

\begin{remark}[A related matrix-group typo]
Section~5.1 says that $\sigma_i(\gamma)$ lies in $\operatorname{SL}_2(\mathbb R)$ while simultaneously imposing $\det(\gamma)\OL=\mathfrak l$.  For a genuine $\mathfrak l$-Hecke matrix one has $\det(\sigma_i(\gamma))=\sigma_i(\pi)$, not $1$.  The M\"obius action is of course well-defined for $\operatorname{GL}_2^+(\mathbb R)$ (or projectively), so this is repairable, but it reinforces that the determinant is the local element $\pi$.
\end{remark}


\section{Classification of the integral norm fiber}
The cyclotomic structure makes the corrected norm equation completely explicit.

\begin{theorem}[The $84$ elements and $42$ traces]\label{thm:42}
Let $\pi$ and $\alpha_0$ be as in \eqref{eq:pi} and \eqref{eq:alpha0}. Then
\begin{equation}\label{eq:normfiber}
 \{x\in\OM:x\bar x=\pi\}
 =\{\zeta_{42}^k\alpha_0:0\le k<42\}
 \sqcup
 \{\zeta_{42}^k\bar\alpha_0:0\le k<42\}.
\end{equation}
In particular the norm fiber contains exactly $84$ elements. Modulo complex conjugation it has $42$ classes, and the corresponding relative traces
\[
 c_k=\zeta_{42}^k\alpha_0+\zeta_{42}^{-k}\bar\alpha_0\in\OL
\]
are pairwise distinct.
\end{theorem}
\begin{proof}
Because $(\alpha_0)=\mathfrak P$ is a prime of $M$ above $\mathfrak l$ and
\[
 (\pi)=\mathfrak P\bar{\mathfrak P},
\]
any integral $x$ with $(x)(\bar x)=(\pi)$ must have $(x)=\mathfrak P$ or $(x)=\bar{\mathfrak P}$. Thus $x=u\alpha_0$ or $x=u\bar\alpha_0$ for a unit $u$ satisfying $u\bar u=1$.

If $u\bar u=1$, then for every complex embedding $\tau:M\hookrightarrow\mathbb C$ one has $|\tau(u)|^2=\tau(u\bar u)=1$.  Since $u$ is an algebraic integer, Kronecker's theorem implies that $u$ is a root of unity.  Conversely every root of unity has relative norm one.  The root-of-unity group of $M=\Ql(\zeta_{42})$ is $\mu_{42}$, so the norm-one unit kernel is exactly $\mu_{42}$, giving \eqref{eq:normfiber}. The two displayed sets are disjoint because they generate the distinct prime ideals $\mathfrak P$ and $\bar{\mathfrak P}$.

Conjugation pairs the $84$ elements. If two conjugacy classes had the same trace, their representatives would have the same trace and the same norm $\pi$, hence would be roots of the same quadratic $X^2-cX+\pi$; they would therefore be equal or conjugate. Thus there are exactly $42$ distinct traces.
\end{proof}

For example,
\[
 c_0=\alpha_0+\bar\alpha_0=t^4-4t^2+t+4.
\]
This theorem replaces an unstructured search for the relative-norm stage by a finite and complete classification for the chosen split prime.


\section{Certified signature enumeration}
For $d=3$ and $d=7$, the same cyclotomic field can be written as $M=L(\sqrt{-d})$. We choose
\[
 \sqrt{-3}=1+2\zeta_{42}^{14}
\]
and
\[
 \sqrt{-7}=1+2(\zeta_{42}^{6}+\zeta_{42}^{12}+\zeta_{42}^{24}),
\]
which square to $-3$ and $-7$, respectively.

For
\[
 \alpha_k=\zeta_{42}^k\alpha_0,
\]
define $a_{k,d},b_{k,d}\in L$ by
\[
 \alpha_k=a_{k,d}+b_{k,d}\sqrt{-d}.
\]
The signs of the six real embeddings of $b_{k,d}$ can be certified exactly by expressing $b_{k,d}$ as a degree-at-most-five polynomial in $t$, isolating the six real roots of (\externalref{packet}{eq:minpoly}) with rational intervals, and evaluating by rational interval arithmetic.

\begin{theorem}[Arithmetic $(3,3)$ signature counts]\label{thm:signature}
With the normalization above:
\begin{enumerate}[label=(\alph*)]
\item for $d=3$, exactly $24$ of the $42$ trace classes have $b_{k,3}$ with three positive and three negative real embeddings;
\item for $d=7$, exactly $4$ of the $42$ trace classes have $b_{k,7}$ with three positive and three negative real embeddings.
\end{enumerate}
A set of representatives is
\[
\begin{split}
 d=3:\quad k\in\{&0,1,4,8,11,12,15,16,17,18,19,20,21,22,25,29,\\
 &32,33,36,37,38,39,40,41\},
\end{split}
\]
and
\[
 d=7:\quad k\in\{2,14,23,35\}.
\]
\end{theorem}

One explicit $d=3$ example is
\[
 b_{0,3}=-\frac13t^5-\frac16t^4+\frac73t^3+t^2-\frac{23}{6}t-\frac43,
\]
whose signs at $k=1,5,11,13,17,19$ are
\[
 (+,-,-,+,+,-).
\]
For $d=7$, one example is
\[
 b_{2,7}=-\frac5{14}t^5+\frac{31}{14}t^3+\frac37t^2-\frac{45}{14}t-\frac{15}{14},
\]
with sign vector
\[
 (+,-,-,+,-,+).
\]

\begin{remark}[What Theorem \ref{thm:signature} does not prove]\label{rem:notgeometry}
The sign vector of $b$ is an arithmetic compatibility condition, not the CM type of a period point. The preprint itself notes that the Hecke endomorphism alone does not determine the CM type. Therefore Theorem \ref{thm:signature} does not prove
\[
 V\cap\mathcal W_K\ne\varnothing.
\]
It proves only that the corrected relative-norm fiber contains integral CM elements whose $b$-coefficients have the desired $3+3$ sign pattern.
\end{remark}


\section{Canonical fixed points and a finite replacement for the lift variable}
The missing elimination of $\delta$ in Section 5.3 can be reorganized using integral conjugacy rather than by ranging directly over $\mathrm{SL}_2(\OL)$. For any corrected trace $c\in\OL$ define the companion matrix
\begin{equation}\label{eq:companion}
 A_c=\begin{pmatrix}0&-\pi\\1&c\end{pmatrix}\in M_2(\OL).
\end{equation}
Then $\det A_c=\pi$ and $\chi_{A_c}(X)=X^2-cX+\pi$.

\begin{proposition}[Canonical Hecke-fixed CM representative]\label{prop:canonical-fixed}
Assume $c^2-4\pi$ is totally negative. For each real embedding $\sigma_i:L\hookrightarrow\R$, put
\begin{equation}\label{eq:wi}
 w_i(c)=\frac{-\sigma_i(c)+\sqrt{\sigma_i(c)^2-4\sigma_i(\pi)}}{2},
\end{equation}
where the square root is chosen with positive imaginary part. Then
\[
 w(c)=(w_1(c),\ldots,w_6(c))\in\mathbb H^6
\]
is fixed componentwise by $A_c$ under the fractional-linear action. Thus $A_c$ supplies an explicit determinant-$\pi$ Hecke-fixed CM representative in the Hilbert modular domain.
\end{proposition}
\begin{proof}
The fixed-point equation for
$\left(\begin{smallmatrix}0&-\pi\\1&c\end{smallmatrix}\right)$ is
$w=-\pi/(w+c)$, equivalently $w^2+cw+\pi=0$. Applying $\sigma_i$ gives \eqref{eq:wi}. Total negativity of the discriminant puts the selected root in $\mathbb H$ at every embedding. The characteristic polynomial generates the same quadratic CM field as a root of $X^2-cX+\pi$.
\end{proof}

For the $42$ traces of Theorem \ref{thm:42}, Proposition \ref{prop:canonical-fixed} therefore gives $42$ explicit canonical fixed-point representatives. They need not be pairwise distinct in the quotient $X_L$, and they do not automatically exhaust all integral conjugacy classes with the same characteristic polynomial.

That last caveat is exactly where the infinite $\delta$-variable belongs. Generalized Latimer--MacDuffee theory replaces conjugacy of integral matrices with a finite module-classification problem for the quadratic order
\[
 R_c=\OL[\alpha],\qquad \alpha^2-c\alpha+\pi=0.
\]
Generalized Latimer--MacDuffee theory identifies the analogous integral problem with module classes over the quadratic order.  For comparison, algorithms for prescribed characteristic and minimal polynomials over $\Z$ are developed in \cite{Marseglia2025}.  In the maximal-order cases needed here, the $\OL$-module classification is handled directly below using class number one. Consequently, a corrected finite search should proceed trace by trace through the finitely many integral conjugacy classes attached to $R_c$, rather than count root choices before the lift has been eliminated.

McMullen's standard modular embedding has the normalization
\begin{equation}\label{eq:mcmullenF}
 F(z)=(z,f_2(z),\ldots,f_6(z))
\end{equation}
with equivariance $f_i(g_1z)=g_i f_i(z)$ \cite[Sec.~6]{McMullen2026}. For the canonical representative \eqref{eq:companion}, the concrete identities
\begin{equation}\label{eq:42test}
 f_i(w_1(c))=w_i(c),\qquad i=2,\ldots,6
\end{equation}
are therefore a sufficient test that the chosen lift $w(c)$ itself lies on the distinguished lift $F(\mathbb H)\subset\mathbb H^6$.

There is an important quotient-level qualification. If $\Gamma_L$ denotes the Hilbert modular group defining $X_L$, then the full inverse image of the McMullen curve in $\mathbb H^6$ is the $\Gamma_L$-saturation
\[
 \bigcup_{\gamma\in\Gamma_L}\gamma\,F(\mathbb H).
\]
Consequently the complete membership criterion for the quotient point $[w(c)]\in X_L$ is
\begin{equation}\label{eq:orbit-test}
 \exists\,\gamma\in\Gamma_L,\ z\in\mathbb H
 \quad\text{such that}\quad
 \gamma\cdot w(c)=F(z).
\end{equation}
Equation \eqref{eq:42test} is the special case $\gamma=1$. For a non-companion integral conjugacy representative one obtains the analogous distinguished-lift test after the corresponding $\mathrm{SL}_2(\OL)$ normalization, but a complete quotient-level search must still account for the orbit condition \eqref{eq:orbit-test}. Thus the arithmetic reductions below make the set of Hecke conjugacy classes finite; they do not by themselves collapse the geometric orbit problem to the single lift $F(\mathbb H)$.


\section{Component normalization and an effective finite orbit search}\label{app:components}
There is a further normalization issue that must be separated from the arithmetic
corrections above.  McMullen's construction first produces the modular embedding in the
disconnected Hilbert modular variety
\[
 Y_L=(\pm\mathbb H)^6/\mathrm{SL}_2(\OL),
\]
not directly in $X_L=\mathbb H^6/\mathrm{SL}_2(\OL)$.  If a Galois-conjugate triangle
has opposite orientation, its component map takes values in $-\mathbb H$ while still
satisfying the equivariance relation \eqref{eq:correct-equiv}; McMullen then uses
commensurability of the components of $Y_L$ and passes to a finite-index subgroup to
obtain a curve on $X_L$ \cite[Prop.~6.1, Cor.~6.2 and proof of Thm.~1.9]{McMullen2026}.

The shortcut in \cite[Sec.~1.4]{Mostaed2026} of replacing a map
$f_i:\mathbb H\to-\mathbb H$ by $-\overline{f_i}$ does not preserve the same target
representation in general.  Let
\[
 J(z)=-\overline z,
 \qquad
 s=\begin{pmatrix}-1&0\\0&1\end{pmatrix}\in\mathrm{PGL}_2(\mathbb R).
\]
For $g\in\mathrm{PSL}_2(\mathbb R)$ one has
\[
 J\circ g\circ J^{-1}=s g s^{-1},
\]
which is not generally equal to $g$.  Hence if
$f_i(g_1z)=g_i f_i(z)$, then $J\circ f_i$ is equivariant for
$g_i\mapsto s g_i s^{-1}$, not for the original $g_i$.  For example, with
$B=\left(\begin{smallmatrix}0&-1\\1&t\end{smallmatrix}\right)$,
\[
 J(Bz)=\frac{1}{\overline z+t},
 \qquad
 B(Jz)=\frac{1}{\overline z-t},
\]
so the two actions differ for $t\ne0$.

Consequently an exact Hecke search on the curve in $X_L$ must specify the finite-cover
and component-changing normalization used to pass from McMullen's $Y_L$-curve to
$X_L$.  This does not alter the abstract local norm calculations of the preceding
sections, but it is part of the data required to identify an arithmetic Hecke class with
a point on the actual curve in $X_L$.

The missing normalization can be written explicitly.  Let
$\varepsilon_i\in\{\pm1\}$ record the component containing $f_i(\mathbb H)$, so that
$f_i(\mathbb H)\subset \varepsilon_i\mathbb H$ and $\varepsilon_1=1$.  By weak
approximation choose $\lambda\in L^\times$ satisfying
$\operatorname{sgn}(\sigma_i(\lambda))=\varepsilon_i$ for all $i$, and put
\[
 h=\begin{pmatrix}\lambda&0\\0&1\end{pmatrix}\in\mathrm{GL}_2(L).
\]
Then the $i$-th real component of $h$ sends $\varepsilon_i\mathbb H$ to $\mathbb H$.
Define
\begin{equation}\label{eq:explicitcovergroup}
 \Gamma_\lambda
 =\Delta_0\cap h^{-1}\mathrm{SL}_2(\OL)h.
\end{equation}
Since $h\mathrm{SL}_2(\OL)h^{-1}$ is commensurable with
$\mathrm{SL}_2(\OL)$, the subgroup $\Gamma_\lambda$ has finite index in $\Delta_0$.
For
\[
 F_\lambda(z)=h\cdot F(z)\in\mathbb H^6
\]
and $g\in\Gamma_\lambda$, equivariance becomes
\[
 F_\lambda(g_1z)=(hgh^{-1})\cdot F_\lambda(z),
 \qquad hgh^{-1}\in\mathrm{SL}_2(\OL).
\]
Thus $F_\lambda$ descends to the finite-cover curve
$\mathbb H/\Gamma_\lambda\to X_L$.  This is the concrete specialization of McMullen's
commensurability argument in Proposition~6.1 and Corollary~6.2.  In particular, after
$\lambda$ is fixed the Hecke-orbit test must be performed against $F_\lambda$, not
against an ad hoc componentwise conjugation of $F$.

For the present $(14,21,42)$ example the component correction is especially simple and
does not itself force a smaller subgroup.

\begin{proposition}[Explicit component orientation and integral transfer]\label{prop:orientationt}
Order the six embeddings by $k=1,5,11,13,17,19$.  The orientations of the six
Galois triangles relative to the identity triangle are
\[
 (\varepsilon_1,\ldots,\varepsilon_6)=(+,+,-,-,-,-).
\]
The unit
\[
 t=2\cos(\pi/21)\in\OL^\times
\]
has the same real sign vector.  Consequently, with
\[
 h_t=\begin{pmatrix}t&0\\0&1\end{pmatrix}\in\mathrm{GL}_2(\OL),
\]
the map $F_t=h_t\cdot F$ takes values in $\mathbb H^6$, and
$h_t\mathrm{SL}_2(\OL)h_t^{-1}=\mathrm{SL}_2(\OL)$.  Thus the component-changing
normalization may be carried out integrally on $X_L$ itself; no additional finite-index
subgroup is required solely for the purpose of changing the signs of the six half-plane
components.
\end{proposition}
\begin{proof}
Use McMullen's printed matrices $A,B,C\in\mathrm{SL}_2(\OL)$ for three consecutively
ordered vertices of the quadrilateral model of $\Delta_0$.  For each of the six real
embeddings indexed by $k=1,5,11,13,17,19$, let $p_A,p_B,p_C\in\mathbb H$ be their
upper-half-plane fixed points.  The oriented Euclidean determinant
\[
 \Omega_k=\det\!\begin{pmatrix}
 \Re(p_B-p_A)&\Im(p_B-p_A)\\
 \Re(p_C-p_A)&\Im(p_C-p_A)
 \end{pmatrix}
\]
has certified rational-interval signs
\[
 \Omega_1<0,\quad \Omega_5<0,\quad
 \Omega_{11}>0,\quad\Omega_{13}>0,\quad\Omega_{17}>0,\quad\Omega_{19}>0.
\]
The six isolating intervals are disjoint and each has an exact sign change for the degree-six
minimal polynomial, so they contain all six real roots, one apiece.  Interval evaluation of
$\Omega_k$ then excludes zero in every case.  Thus, relative to $k=1$, the orientation
vector is $(+,+,-,-,-,-)$.  The same sign
vector for $t$ is immediate from
$\sigma_k(t)=2\cos(k\pi/21)$.  Finally, the minimal polynomial
(\externalref{packet}{eq:minpoly}) has constant term $1$, so
$N_{L/\Ql}(t)=1$ and $t\in\OL^\times$.  Hence $h_t\in\mathrm{GL}_2(\OL)$ and
conjugation by $h_t$ preserves $\mathrm{SL}_2(\OL)$.
\end{proof}

A torsion-free or level cover may still be imposed for fine-moduli purposes, but the
component transfer itself is already integral.  This removes one previously unspecified
normalization from the Hecke comparison: in what follows the normalized lift can be
taken to be $F_t$.

Once that normalization is fixed, compactness gives an effective replacement for the
apparently infinite orbit search in \eqref{eq:orbit-test}.

\begin{proposition}[Effective finiteness of the orbit test]\label{prop:effectiveorbit}
Let $G\subset\mathrm{SL}_2(L)$ be an arithmetic subgroup commensurable with
$\mathrm{SL}_2(\OL)$ acting diagonally on $\mathbb H^6$, let
$F:\mathbb H\to\mathbb H^6$ be a $\Gamma$-equivariant modular embedding with
$\Gamma$ cocompact, and let $D\subset\mathbb H$ be a compact fundamental polygon for
$\Gamma$.  Put $K=F(D)$.  For every fixed $w=(w_1,\ldots,w_6)\in\mathbb H^6$, the set
\[
 \{\gamma\in G:\gamma w\in K\}
\]
is finite.  Moreover, after choosing an integral lattice containing $G$, it is
effectively enumerable from certified lower bounds for the coordinate imaginary parts
on $K$ and upper bounds for the coordinate absolute values on $K$.
\end{proposition}
\begin{proof}
Write $w_i=x_i+i y_i$ with $y_i>0$, and
$\gamma_i=\left(\begin{smallmatrix}a_i&b_i\\c_i&d_i\end{smallmatrix}\right)$ for the
$i$-th real embedding of $\gamma$.  Since $K$ is compact in $\mathbb H^6$, there are
constants
\[
 m_i=\inf_{z\in K}\operatorname{Im}(z_i)>0,
 \qquad
 U_i=\sup_{z\in K}|z_i|<\infty.
\]
If $u=\gamma w\in K$, the Möbius imaginary-part identity gives
\[
 \operatorname{Im}(u_i)=\frac{y_i}{|c_iw_i+d_i|^2}\ge m_i,
\]
hence
\[
 |c_iw_i+d_i|\le B_i:=\sqrt{y_i/m_i}.
\]
Taking imaginary parts yields
\[
 |c_i|\le\frac{B_i}{y_i}=\frac1{\sqrt{m_i y_i}},
\]
and therefore
\[
 |d_i|\le B_i+|c_i|\,|w_i|.
\]
Likewise $a_iw_i+b_i=u_i(c_iw_i+d_i)$, so
\[
 |a_iw_i+b_i|\le U_iB_i.
\]
Taking imaginary parts once more and then using the triangle inequality gives
\[
 |a_i|\le\frac{U_iB_i}{y_i},
 \qquad
 |b_i|\le U_iB_i+|a_i|\,|w_i|.
\]
Thus every real embedding of every matrix entry is bounded by an explicit constant.
The Minkowski embedding of the chosen coefficient lattice is discrete in a finite-dimensional
real vector space, so only finitely many lattice elements satisfy these simultaneous
archimedean bounds.  Enumerating them and imposing determinant one gives the claimed
effective finite set.
\end{proof}

\begin{corollary}[Finite geometric search for each CM candidate]\label{cor:finitegeometric}
After the explicit component normalization of Proposition \ref{prop:orientationt} (and any independently chosen level cover), the quotient-membership test for each of the $12$ strongest maximal-order
$d=3$ candidates reduces to a finite computation.  In particular, the infinite union in
\eqref{eq:orbit-test} does not require an unbounded search through
$\mathrm{SL}_2(\OL)$: only matrices satisfying the archimedean bounds of
Proposition \ref{prop:effectiveorbit} can occur.
\end{corollary}

The generic finite-search argument requires a specific finite-index model
of the curve in $X_L$, a compact polygon $D$, and certified bounds
$m_i,U_i$ for its normalized component maps. The corrected Schwarz equations
in Proposition~\ref{prop:sixHG} provide one route to such bounds.
For the fixed packet considered here, the explicit two-region construction
of Lemma~\externalref{packet}{lem:twochambers} and the exact certificates assembled in
Theorem~\externalref{packet}{thm:fullpacketgeometry} complete the geometric decisions.


\section{Order-discriminant stratification of the $42$ traces}\label{app:orders}
There is a further exact reduction at $43$. The field discriminants are
\[
 \operatorname{disc}(L)=453789=3^3 7^5,
 \qquad
 \operatorname{disc}(M)=205924456521=3^6 7^{10}.
\]
Hence
\[
 |\operatorname{disc}(M)|=\operatorname{disc}(L)^2.
\]
By the discriminant tower formula, the relative discriminant ideal $\mathfrak D_{M/L}$ has norm one and is therefore the unit ideal. In particular, $M/L$ is unramified at all finite primes.

For $\alpha_k=\zeta^k\alpha_0$ and $c_k=\alpha_k+\bar\alpha_k$, the relative discriminant of the basis $(1,\alpha_k)$ is
\begin{equation}\label{eq:Deltak}
 \Delta_k=(\alpha_k-\bar\alpha_k)^2=c_k^2-4\pi.
\end{equation}
Since the maximal relative discriminant is the unit ideal, $\OL[\alpha_k]=\OM$ exactly when $(\Delta_k)=\OL$.

\begin{theorem}[Maximal-order trace classes at $43$]\label{thm:maxorder}
Among the $42$ trace classes of Theorem \ref{thm:42}, exactly $16$ generate the maximal CM order $\OM$. With the indexing $\alpha_k=\zeta^k\alpha_0$, they are
\[
 k\in\{1,7,11,13,16,17,19,20,22,28,32,34,37,38,40,41\}.
\]
The absolute norms of the relative order discriminants across all $42$ classes belong to
\[
 \{1,49,1681,1849,2401,6889,16129\}
 =\{1,7^2,41^2,43^2,49^2,83^2,127^2\}.
\]
Exactly $12$ of the $16$ maximal-order classes are also among the $24$ $d=3$ arithmetic $(3,3)$ classes, namely
\[
 \{1,11,16,17,19,20,22,32,37,38,40,41\}.
\]
None of the four $d=7$ arithmetic $(3,3)$ classes is maximal-order.
\end{theorem}
\begin{proof}
Equation \eqref{eq:Deltak} is evaluated exactly in the power basis of $L$, and its absolute norm is the resultant with the minimal polynomial (\externalref{packet}{eq:minpoly}). The resulting norm list is printed in the computational certificate. Because $\mathfrak D_{M/L}=\OL$, norm one is equivalent to equality of the generated order with $\OM$. Intersecting the resulting $16$ indices with the certified lists in Theorem \ref{thm:signature} gives the final two assertions.
\end{proof}

\begin{remark}
The maximal-order cases are particularly attractive for the conjugacy repair. The cyclotomic field $M=\Ql(\zeta_{42})=\Ql(\zeta_{21})$ has class number one by Masley's Main Theorem \cite[p.~179]{Masley1976}. This removes the ordinary ideal-class obstruction in $\OM$; however, one must still track the precise $\OL$-module, determinant, and $\mathrm{SL}_2$ versus $\mathrm{GL}_2$ equivalence conditions before concluding uniqueness of an admissible Hecke conjugacy class.
\end{remark}

\begin{proposition}[All-trace reduction by embedded order]\label{prop:alltraceorders}
The $42$ traces generate exactly nine embedded orders
$R_k=\OL[\alpha_k]$ in the fixed field $M$. Their conductors in $\OM$
are the extensions of the following ideals of $\OL$:
\begin{center}
\begin{tabular}{ll}
\toprule
Conductor in $\OL$ & Trace indices $k$\\
\midrule
$\OL$ & $1,7,11,13,16,17,19,20,22,28,32,34,37,38,40,41$\\
$(7,t-1)$ & $0,3,6,9,15,18,21,24,27,30,36,39$\\
$(7,t-1)^2$ & $12,33$\\
$(41,t-12)$ & $14,35$\\
$(41,t-22)$ & $5,26$\\
$(43,t-20)$ & $4,25$\\
$(43,t-31)$ & $8,29$\\
$(83,t-36)$ & $2,23$\\
$(127,t-94)$ & $10,31$\\
\bottomrule
\end{tabular}
\end{center}
Within each row, an integral invertible affine change of generator gives
an $\SL_2(\OL)$-equivariant correspondence of all integral actions,
preserving their upper-half-plane fixed tuples. For every trace, the
maximal-multiplier population is represented by the same sixteen maximal
points in Theorem~\ref{thm:d7ambient}.
\end{proposition}
\begin{proof}
Write $\alpha_k=a_k+b_k\zeta$ in $\OM=\OL+\OL\zeta$.
Then $R_k=\OL+b_k\OM$, and its conductor is $b_k\OM$:
if $A+B\zeta$ and its product with $\zeta$ lie in $R_k$, then both
$B$ and $A+tB$ are divisible by $b_k$, so $A$ is too.
For each prime-conductor row, the two exact implementations listed in
Section~\ref{app:computer} verify $|N_{L/\mathbf Q}(b_k)|=p$ and
$b_k(r)=0\pmod p$ for the displayed root $r$.
Inclusion and equal index give $(b_k)=(p,t-r)$.
For the square-conductor row, $b_k/(t-1)^2$ and its inverse are integral
units, and $|N_{L/\mathbf Q}(t-1)|=7$.
All within-row ratios $e=b_l/b_k$ are checked to be integral units.
Thus $\alpha_l=s+e\alpha_k$, with $s=a_l-ea_k\in\OL$, and
$T\mapsto sI+eT$ is the required correspondence. The fixed-point
equation is multiplied by the nonzero scalar $e$ at each embedding.

For maximal multiplier, multiplication by $\alpha_k$ in an orientation
sector $\delta\in\OL^\times$ has matrix
\[
 \begin{pmatrix}a_k&-\delta b_k\\ b_k/\delta&a_k+tb_k\end{pmatrix}.
\]
Here $b_k\ne0$, since otherwise $\pi=a_k^2$ would have square absolute
norm. Its upper-half-plane tuple has
\[
 x=-\delta t/2,\qquad y^2=\delta^2(4-t^2)/4,
 \qquad |w|^2=\delta^2,
\]
independently of $k$. The maximal Picard group is trivial, and
Theorem~\ref{thm:unitnorm} gives the same $32$ determinant sectors for
every trace. The root identifications in Theorem~\ref{thm:d7ambient}
reduce them to the same sixteen representatives.
\end{proof}
\begin{remark}
The two conductor ideals of norm $41$ are different embedded ideals.
Also the basis change from $(1,\alpha_k)$ to $(1,\alpha_l)$ has
determinant $e$, so companion orientation labels need not be unchanged.
This proposition gives no count of all proper modules and no new curve
decision; those require the separate class and orbit analyses.
\end{remark}

\begin{theorem}[Exact nonmaximal $d=7$ conjugacy counts]\label{thm:d7nonmax}
For the four $d=7$ arithmetic $(3,3)$ trace classes
\[
 k\in\{2,14,23,35\},
\]
put $R_k=\OL[\alpha_k]\subset\OM$.  Their conductors are
\[
 \mathfrak f_2=\mathfrak f_{23}=(83,t-36),\qquad
 \mathfrak f_{14}=\mathfrak f_{35}=(41,t-12).
\]
The two conductor primes have residue degree one in $L$ and are inert in the quadratic extension $M/L=L(\sqrt{-7})/L$.  The proper ideal class groups and full ideal-class monoids are
\[
\begin{array}{c|c|c|c}
 k & N\mathfrak f_k & |\operatorname{Pic}(R_k)| &
 |\operatorname{ICM}(R_k)|\\ \hline
 14,35&41&1&2\\
 2,23&83&2&3.
\end{array}
\]
Within any fixed totally-positive determinant-orientation sector, refining the corresponding $\mathrm{GL}_2^+(\OL)$ similarity classes to determinant one gives
\[
 3\quad(k=14,35),\qquad 5\quad(k=2,23)
\]
classes.  Across all determinant-sign sectors, the full numbers of $\mathrm{SL}_2(\OL)$ similarity classes are
\[
 96\quad(k=14,35),\qquad 160\quad(k=2,23).
\]
Thus the nonmaximal $d=7$ conjugacy problem is completely finite, but the full oriented enumeration is substantially larger than the single totally-positive sector.
\end{theorem}
\begin{proof}
Since the maximal relative discriminant of $M/L$ is the unit ideal, the order-discriminant formula gives
\[
 \operatorname{disc}(R_k/\OL)=\mathfrak f_k^2.
\]
Exact reduction of $\Delta_k=c_k^2-4\pi$ gives
\[
 N(\Delta_2)=N(\Delta_{23})=83^2,
 \qquad
 N(\Delta_{14})=N(\Delta_{35})=41^2,
\]
and modulo the completely split rational primes one finds that the unique vanishing factors are $t\equiv36\pmod{83}$ and $t\equiv12\pmod{41}$, respectively.  Thus the displayed conductor ideals are exact.  Moreover
\[
 \left(\frac{-7}{41}\right)=\left(\frac{-7}{83}\right)=-1,
\]
so the conductor primes are inert in $M/L$.  Therefore locally
\[
 (\OM/\mathfrak f_k\OM)^\times/(\OL/\mathfrak f_k)^\times
 \cong \mathbb F_{p^2}^\times/\mathbb F_p^\times
\]
is cyclic of order $p+1$.

Let $z=\zeta_{42}$, $\rho=z^2=\zeta_{21}$, and $\varepsilon=1-\rho$.  Because $\Phi_{21}(1)=1$, $\varepsilon$ is a global unit.  At either conductor prime $p\equiv-1\pmod{42}$, Frobenius is complex conjugation and
\[
 \varepsilon^{p-1}
 =\frac{1-z^{-2}}{1-z^2}
 =-z^{-2}=z^{19}.
\]
Since $19$ is prime to $42$, the image of $\varepsilon$ in the local quotient has order $42$.
We must also show that no omitted global unit enlarges this image.  For any
$\eta\in E_M$, Kronecker's theorem gives
$\eta/\bar\eta\in\mu(M)$.  The map
\[
 E_M/(\mu(M)E_L)\longrightarrow \mu(M)/\mu(M)^2,
 \qquad \eta\longmapsto\eta/\bar\eta,
\]
is injective: if $\eta/\bar\eta=\omega^2$, then $\eta/\omega\in E_L$.
The target has order two, while
$\varepsilon/\bar\varepsilon=-z^2=z^{23}$ is not a square in $\mu_{42}$.
Consequently
\[
 E_M=\langle\varepsilon\rangle\mu_{42}E_L.
\]
Since $[z]=[\varepsilon]^{22}$ locally, the global-unit image is exactly the
order-$42$ subgroup generated by $[\varepsilon]$.  At the conductor-$83$
prime the element $1+3z$ has exact quotient order $84$, so it represents the
other coset.  The standard order class exact sequence and $h(M)=1$ now give
\[
 |\operatorname{Pic}(R_{14})|=|\operatorname{Pic}(R_{35})|=42/42=1,
\]
\[
 |\operatorname{Pic}(R_2)|=|\operatorname{Pic}(R_{23})|=84/42=2.
\]
For clarity, the ideal-class-monoid completeness can be seen directly in this
prime-conductor case.  Locally write $A=(\OL)_{\mathfrak f}$,
$B=(\OM)_{\mathfrak f}$, and $R=A+\mathfrak fB$.  After scaling a rank-one
torsion-free $R$-lattice $I$, one has $R\subset I\subset B$.  Since $B/R$ is
one-dimensional over the residue field of $A$, necessarily $I=R$ or $I=B$.
Away from $\mathfrak f$ the order is maximal.  Thus the multiplier order is
either $R_k$ or $\OM$, and within either multiplier order the classes form its
Picard group.  Equivalently, this is the prime-conductor Bass-order
decomposition
\[
 \operatorname{ICM}(R_k)
 =\operatorname{Pic}(R_k)\sqcup\operatorname{Pic}(\OM).
\]
Since $h(M)=1$, the displayed ideal-class-monoid counts follow.

For the determinant-one refinement, the integral centralizer of an $R_k$-multiplier class is $R_k^\times$.  In the same local quotient one has $[z]=[\varepsilon]^{22}$.  Hence a unit $\varepsilon^n z^m u$, $u\in E_L$, lies in $R_k^\times$ only if
\[
 n+22m\equiv0\pmod{42},
\]
which forces $n$ even.  Conversely every square of an $E_L$-unit is a norm
from $R_k^\times$.  Hence the norm calculation sharpens to
\[
 N_{M/L}(R_k^\times)=E_L^2.
\]
Because $[E_L^+:E_L^2]=2$, every $R_k$-multiplier class splits into two determinant-one classes inside a fixed totally-positive determinant sector, whereas the $\OM$-multiplier class has centralizer $\OM^\times$ and does not split there.  This gives $2|\operatorname{Pic}(R_k)|+1$, namely $3$ for conductor $41$ and $5$ for conductor $83$, in such a sector.

For the full $\mathrm{SL}_2(\OL)$ enumeration one must also retain determinant signs.  The maximal-overorder class splits according to
\[
 E_L/N(\OM^\times)=E_L/E_L^+,
\]
of order $32$, while each $R_k$-multiplier class splits according to
\[
 E_L/N(R_k^\times)=E_L/E_L^2,
\]
of order $64$.  Hence the total counts are $64|\operatorname{Pic}(R_k)|+32$: $96$ for conductor $41$ and $160$ for conductor $83$.  The primary and independent exact certificates construct every module template, all determinant-sign twists, and all $512$ trace-labelled representatives.
\end{proof}

\begin{theorem}[The nonmaximal $d=3$ classes]\label{thm:d3classes}
For the twelve selected arithmetic-$(3,3)$ traces
\[
 k\in\{0,4,8,12,15,18,21,25,29,33,36,39\},
\]
there are exactly $1280$ trace-labelled $\SL_2(\OL)$ matrix classes,
or $40$ in one chosen totally positive determinant sector per module class.
These are matrix-class counts, not counts of distinct ambient or curve points.
The multiplier orders have conductors
$\OL$, $(7,t-1)$, $(7,t-1)^2$, $(43,t-20)$, or $(43,t-31)$,
and each of these orders has trivial Picard group.
\end{theorem}
\begin{proof}
Proposition~\ref{prop:alltraceorders} supplies the conductors.
Ten labels have prime conductor and two, $12,33$, have the square conductor.
Every relevant prime splits in $M/L$. In the conductor unit quotient
\[
 Q_{\mathfrak f}=(\OM/\mathfrak f\OM)^\times/
                    (\OL/\mathfrak f)^\times,
\]
use the ratio of the two split coordinates. With
$\varepsilon=1-z^2$ as above, the following exact table gives generators:
\begin{center}
\begin{tabular}{llll}
\toprule
$\mathfrak f$ & Residue ring & $z$ & $[\varepsilon]$, order\\
\midrule
$(7,t-1)$ & $\mathbf F_7$ & $3$ & $5$, $6$\\
$(7,t-1)^2$ & $\mathbf F_7[v]/(v^2)$ & $3+2v$ & $5+2v$, $42$\\
$(43,t-20)$ & $\mathbf F_{43}$ & $30$ & $3$, $42$\\
$(43,t-31)$ & $\mathbf F_{43}$ & $5$ & $18$, $42$\\
\bottomrule
\end{tabular}
\end{center}
In the second row $t=1+v$: the identity
$f(1+v)=7(1-v)\pmod{v^2}$ proves the displayed residue-ring description.
In each row the class of $\varepsilon$ generates the entire quotient.
The conductor exact sequence and $h(M)=1$ therefore give trivial Picard
groups for these four proper orders. Moreover $[z]=[\varepsilon]^{22}$.
The complete relative unit generation
$E_M=\langle\varepsilon\rangle\mu_{42}E_L$ proved above shows that
$\varepsilon^n z^m w$, $w\in E_L$, is a unit of the proper order
precisely when $n+22m\equiv0\pmod d$, where $d=6$ or $42$.
As $\gcd(22,d)=2$, its exponent $n$ is even. Consequently the proper-order
norm image is exactly $E_L^2$; the maximal-order image is $E_L^+$.

For completeness of the module list, let $D$ be the completed local DVR,
$B=D\times D$, and $R_e=D+\varpi^eB$, with $e=1$ or $2$.
Scale any full rank-one $R_e$ lattice $I$ so that $BI=B$.
Both coordinate projections of $I$ are surjective. Its image in the residue
plane thus contains a vector with both coordinates nonzero, so $I$ contains
a unit of $B$. Scaling by its inverse gives $R_e\subseteq I\subseteq B$.
The quotient $B/R_e\simeq D/(\varpi^e)$ is cyclic; its submodules have
preimages exactly $R_j=D+\varpi^jB$, $0\le j\le e$.
Their multiplier rings are distinct. Away from the conductor every lattice
is locally principal over the maximal order. Globally every lattice is
therefore invertible over its own multiplier order, and its class belongs
to that order's Picard group. This proves there are exactly $e+1$ module
classes, represented by the overorders themselves, including the intermediate
conductor-7 order in the square-conductor case.

Since $h(L)=1$, each lattice is free of rank two over $\OL$.
The characteristic polynomial is irreducible over $L$, so integral matrix
classes correspond to these rank-one $R_k$ modules.
The determinant of a centralizer unit is its relative norm; a module with
multiplier $S$ therefore contributes $E_L/N(S^\times)$ determinant-one
classes. This gives $32$ for maximal multiplier and $64$ for each proper
multiplier, or $1$ and $2$ in a fixed positive determinant sector.
The totals are $10(32+64)+2(32+2\cdot64)=1280$ and
$10(1+2)+2(1+2\cdot2)=40$.
The two exact implementations reconstruct all representatives and check
their traces and determinants; completeness and nonconjugacy follow from
the module and determinant quotients just proved.
\end{proof}

\begin{theorem}[The $112$ ambient nonmaximal $d=7$ CM points]\label{thm:d7ambient}
Let $\mathcal X_{41}$ and $\mathcal X_{83}$ be the sets of
$\mathrm{SL}_2(\OL)$-quotient upper fixed points arising from the trace pairs
$\{14,35\}$ and $\{2,23\}$, respectively.  Then
\[
 |\mathcal X_{41}|=48,\qquad |\mathcal X_{83}|=80,
 \qquad |\mathcal X_{41}\cap\mathcal X_{83}|=16,
\]
and hence
\[
 |\mathcal X_{41}\cup\mathcal X_{83}|=112.
\]
The union has the disjoint multiplier-order decomposition
\[
 32\text{ proper-}R_{41}
 \;\sqcup\;
 64\text{ proper-}R_{83}
 \;\sqcup\;
 16\text{ maximal-}\OM.
\]
Thus the $512$ objects of Theorem~\ref{thm:d7nonmax} are trace-labelled
oriented matrix classes, not $512$ distinct quotient points.

The normalized $A,B,C$ vertices are three of the $16$ maximal-multiplier
points.  Exact determinant-orientation fingerprints identify two further
maximal points with the companion pairs $\{1,22\}$ and $\{17,38\}$ excluded
in Proposition~\ref{prop:residualexclusions}. These vertex and residual
results alone leave $107$ ambient membership decisions: $11$
maximal-multiplier and $96$ proper-order points. The further exclusions below complete these historical v45 membership
decisions in Corollary~\ref{cor:d7geometrycomplete}.
\end{theorem}
\begin{proof}
Fix one trace $c=c_k$.  For every elliptic matrix $T$, the matrix
$\bar T=cI-T$ has the same fixed-point equation.  The involution
$T\mapsto\bar T$ is free on $\mathrm{SL}_2(\OL)$ similarity classes.  Indeed,
at every real embedding a conjugator from $T$ to $\bar T$ must exchange the
upper and lower eigenlines, and therefore has negative determinant.  More
algebraically, every anti-intertwiner is complex conjugation followed by an
element $m$ of the CM centralizer, so its determinant has the totally negative
form $-N_{M/L}(m)$.  It cannot have determinant one.

Conversely, if two classes give the same quotient point, conjugate them so
that they fix the same upper tuple.  Their real coefficients make them fix the
same lower tuple as well, so both lie in the same quadratic stabilizer algebra.
The polynomial $X^2-cX+\pi$ has only its two roots there.  Hence every fixed
point fiber has exactly the two classes $T$ and $cI-T$.  The $96$ and $160$
classes therefore give $48$ and $80$ points for a fixed conductor trace.

The exact identities
\[
 \alpha_{23}=-\alpha_2,\qquad \alpha_{35}=-\alpha_{14}
\]
show that the two traces for each conductor have the same point set.  If
\[
 \gamma_{41}=-1-2t+t^2-t^3-t^4,
 \qquad \gamma_{83}=-2-2t+2t^2-t^4,
\]
then exact reduction gives
\[
 (-174-718t-267t^2+733t^3+131t^4-167t^5)\gamma_{41}
 +30\gamma_{83}=1.
\]
Thus the conductor ideals are comaximal.  A point supporting both proper
orders has multiplier ring containing the order they generate, namely $\OM$.
Conversely every maximal-multiplier point supports all four $\alpha_k$.
The intersection is therefore exactly the maximal stratum, whose $32$ matrix
classes give $16$ points.  Halving within each multiplier stratum yields
\[
64/2=32,\qquad 128/2=64,\qquad 32/2=16,
\]
and proves the $112$-point decomposition.  Equivalently, each proper point
occurs four times in the trace-labelled manifest and each maximal point eight
times, giving $32\cdot4+64\cdot4+16\cdot8=512$.

Finally, the integral generators $Z_A,B_t,C_t$ all generate $\OM$, so direct
and conjugate evaluation of every $\alpha_k$ gives the complete two-class
fiber over each of the three distinct vertices.  In the exact unit-sign basis
$(\xi_5,\xi_{11},\xi_{13},\xi_{17},-1,4-t^2)$ their maximal-point masks are
$7,15,0$.  The companion bridges for $\{1,22\}$ and $\{17,38\}$ have masks
$13$ and $2$.  These five masks are distinct, leaving $11+96=107$ ambient
points.  Both exact implementations verify every representative, root
pairing, and mask.
\end{proof}


